\begin{proof}[Proof of \cref{obj:lem:false-light-envelope}]
All expectations below are under the iid infinite product law generating the observation sequence, with each statistic evaluated on its first \(n\) coordinates.  Let
\[
L=\log(en),\qquad \ell=\log n,
\qquad m_j=|I_j|\ (j=0,1),
\qquad B=B(n),\qquad M=M(n),
\qquad G=6^{2M}.
\]
The deterministic split construction of \cref{obj:def:sample-splits} gives
\[
m_0=\lfloor n/2\rfloor,
\qquad
m_1=n-\lfloor n/2\rfloor .
\]
By \cref{obj:lem:calibration-consequences},
\[
0<\ell,
\qquad \ell\le L,
\qquad L\le 2\ell,
\qquad m_1>0,
\qquad n\le 2m_1,
\qquad M\le L,
\]
\[
B\le\frac{8192L}{n},
\qquad G^2L^6\le n .
\]
% lean: large_calibration_bounds
In particular \(n>0\).  Since \(B=4096L/m_1\), \(L>0\), and \(m_1>0\), also \(B>0\), and hence
\[
B^2\le \frac{8192^2L^2}{n^2}.
\]

Because \(n>0\), \(L=1+\ell\).  Since \(L\ge240\),
\[
\exp(-49L)\le \exp(-L)=\frac1{en}\le \frac1n .
\]
Moreover \(G\ge1\), so \(L^4\le L^6\le G^2L^6\le n\).  Together with \(\ell\le L\), this gives
\[
L^2\ell^2\le L^4\le n.
\]
Therefore
\[
\frac{d^2L^2}{n^3}
\le
\frac{d^2}{n^2\ell^2}
\le
\frac1n+\frac{d^2}{n^2\ell^2}.
\]
% lean: selectedFalseLightError_rate

Let
\[
F=\{k:p_k>B/4\},
\qquad
S_k=\mathbf 1\{k\in\widehat{\mathcal L}_n\},
\qquad
e_k=S_k(\widehat P_k-\phi(q_k)).
\]
By the definition of the false-light error, the selected false-light error is \(\sum_{k\in F}e_k\).  For each sample, Cauchy's inequality gives
\[
\left(\sum_{k\in F}e_k\right)^2
\le
|F|\sum_{k\in F}e_k^2 .
\]

Fix \(k\in F\).  The cutoff defining \(N_0\) in \cref{obj:def:hybrid-estimator-handle} and the assumption \(n\ge N_0\) give
\[
\alpha_0L\ge2,
\qquad
4M^2\le m_1,
\qquad
\frac{4M}{m_1}\le \frac{3B}{4}.
\]
Since \(M\ge2\), the bound \(4M^2\le m_1\) gives \(16\le m_1\).  With the split identities above, this implies \(m_0>0\) and \(m_1\le2m_0\).

Write
\[
G_M(x)=\sum_{j=0}^{M-2}g_{M,j}x^j,
\qquad
 g_{M,j}=(-1)^j
 \frac{2^{2j+3}}{M(M+j+2)}\binom{M+j+2}{2j+4}.
\]
Let \(\mathbf e_{by}\) denote the unit multi-index at the four-cell coordinate \((b,y)\in\{0,1\}^2\).  In the next display, \(ay\) is a single four-cell label, say \(ay=(b,y')\), and \(\mathbf e_{ay}\) is the corresponding unit multi-index.  For \(a\in\{0,1\}\), \(ay\in\{0,1\}^2\), and \(0\le t\le j\le M-2\), set
\[
r(a,ay,j,t)
=
\mathbf e_{ay}+\mathbf e_{a1}+t\mathbf e_{a0}+(j-t)\mathbf e_{a1},
\qquad |r(a,ay,j,t)|=j+2.
\]
For a multi-index \(r\), the factorial-moment monomial is
\[
\mathfrak m_{k,r}
=
\frac{\prod_{(b,y)\in\{0,1\}^2}
\bigl(N^{(1)}_{bky}\bigr)_{r_{by}}}
{(m_1)_{|r|}},
\]
where \((z)_s\) denotes the falling factorial.  The \(a\)-arm sparse factorial-moment contribution is
\[
\widehat P_{k,a}^{\mathrm{sp}}
=
\sum_{j=0}^{M-2}\sum_{t=0}^j\sum_{ay\in\{0,1\}^2}
B^{-1}g_{M,j}B^{-j}\binom jt\,
\mathfrak m_{k,r(a,ay,j,t)}.
\]
By \cref{obj:def:hybrid-estimator-handle}, \(\widehat P_k=\sum_{|r|\le M}c_r\,\mathfrak m_{k,r}\) in collected coefficient form, where each \(c_r\) is obtained by expanding the two arms of \(P_{M,B}\) through the coefficients \(g_{M,j}\) and the binomial expansion of the arm masses, and summing the resulting sparse coefficients \(c_u=B^{-1}g_{M,j}B^{-j}\binom jt\) over the sparse indices \(u=(j,t,ay)\) with \(r_u=r\).  Since \(r\mapsto\mathfrak m_{k,r}\) enters linearly in each coefficient, regrouping this finite double sum arm by arm yields
\[
\widehat P_k=\widehat P_{k,1}^{\mathrm{sp}}-
\widehat P_{k,0}^{\mathrm{sp}} .
\]

For a nonnegative four-cell vector \(v\), define the absolute sparse arm envelope
\[
\mathcal A_a(v)
=
\sum_{j=0}^{M-2}\sum_{t=0}^j\sum_{ay\in\{0,1\}^2}
\left|B^{-1}g_{M,j}B^{-j}\binom jt\right|
 v^{r(a,ay,j,t)},
\qquad
v^r=\prod_{(b,y)\in\{0,1\}^2}v_{by}^{r_{by}}.
\]
Set
\[
v^{(k)}_{ay}=q_{k,ay}+\frac{M}{m_1}.
\]
This vector is nonnegative.  Fix two sparse indices \(u=(j,t,ay)\) and \(u'=(j',t',ay')\) with \(0\le t\le j\le M-2\) and \(0\le t'\le j'\le M-2\).  Both multi-indices \(r_u,r_{u'}\) have degree at most \(M\), so \(4\max\{|r_u|,|r_{u'}|\}^2\le4M^2\le m_1\).  Apply \cref{obj:lem:factorial-product-moment} to the degree-ordered pair.  Its unshifted monomial factor is bounded coordinatewise by \((v^{(k)})^{r}\), since \(q_{k,cy}\le q_{k,cy}+M/m_1\).  Its shifted monomial factor is also bounded coordinatewise by the corresponding power of \(v^{(k)}\): if \(s\) is the larger-degree multi-index in the ordered pair, then \(|s|\le M\), and hence
\[
q_{k,cy}+\frac{|s|}{m_1}
\le
q_{k,cy}+\frac{M}{m_1}
=v^{(k)}_{cy}
\qquad\text{for every coordinate }cy.
\]
Using these two coordinatewise bounds, and swapping \(u,u'\) before applying \cref{obj:lem:factorial-product-moment} when the opposite degree order holds, gives
\[
E\!\left[
\bigl(c_u\mathfrak m_{k,r_u}\bigr)
\bigl(c_{u'}\mathfrak m_{k,r_{u'}}\bigr)
\right]
\le
\mathrm e\,
\left|c_u\right|(v^{(k)})^{r_u}
\left|c_{u'}\right|(v^{(k)})^{r_{u'}},
\]
where \(c_u=B^{-1}g_{M,j}B^{-j}\binom jt\) and \(r_u=r(a,ay,j,t)\).  Expanding the square of \(\widehat P_{k,a}^{\mathrm{sp}}\), summing this bound over all ordered pairs of sparse indices, and collecting the two sums into the envelope gives
\[
E\!\left[(\widehat P_{k,a}^{\mathrm{sp}})^2\right]
\le
\mathrm e\,\mathcal A_a(v^{(k)})^2,
\qquad a=0,1.
\]

The shifted vector has total mass
\[
\sum_{ay}v^{(k)}_{ay}
=p_k+\frac{4M}{m_1}
\le p_k+\frac{3B}{4}
\le B\left(1+\frac{p_k}{B}\right).
\]
Let \(R=\sum_{ay}v_{ay}\), \(A_a(v)=v_{a0}+v_{a1}\), and
\[
g_M^+(z)=\sum_{j=0}^{M-2}|g_{M,j}|z^j.
\]
For every nonnegative \(v\) with \(R\le Bw\) and \(w\ge1\), the binomial summation in the sparse envelope gives
\[
\mathcal A_a(v)
=
B^{-1}R\,v_{a1}\,g_M^+\!\left(\frac{A_a(v)}B\right),
\]
and, since \(0\le A_a(v)/B\le R/B\le w\) with \(w\ge1\), \cref{obj:lem:coefficient-sum-certificate} gives
\[
g_M^+\!\left(\frac{A_a(v)}B\right)
\le 6^M\max\left\{1,\left(\frac{A_a(v)}B\right)^{M-2}\right\}
\le 6^M w^{M-2}.
\]
Since \(v_{a1}\le R\le Bw\), these displays imply
\[
\mathcal A_a(v)
\le B6^M w^M.
\]
Applying this with \(v=v^{(k)}\) and \(w=1+p_k/B\) yields
\[
\mathcal A_a(v^{(k)})
\le
B6^M\left(1+\frac{p_k}{B}\right)^M,
\qquad a=0,1.
\]
Since \(\widehat P_k\) is the difference of the two sparse-arm contributions, \((x-y)^2\le2x^2+2y^2\) yields
\[
E[\widehat P_k^2]
\le
4\mathrm e\,
\left\{
B6^M\left(1+\frac{p_k}{B}\right)^M
\right\}^2 .
\]
Because \(k\in F\) and \(B>0\), \(p_k>0\).  The overlap condition and \cref{obj:def:cell-functional} give \(|\phi(q_k)|\le p_k\).  Since \(p_k/B\ge0\) and \(M\ge2\),
\[
\frac{p_k}{B}
\le
6^M\left(1+\frac{p_k}{B}\right)^M,
\]
and therefore
\[
\phi(q_k)^2
\le
\left\{
B6^M\left(1+\frac{p_k}{B}\right)^M
\right\}^2 .
\]
Using \((x-t)^2\le2x^2+2t^2\),
\[
E\!\left[(\widehat P_k-\phi(q_k))^2\right]
\le
8(\mathrm e+1)B^2
\left\{6^M\left(1+\frac{p_k}{B}\right)^M\right\}^2 .
\]

The indicator \(S_k\) is determined by the pilot split, while \(\widehat P_k-\phi(q_k)\) is computed from the estimation split and a fixed target.  The split factorization gives
\[
E\!\left[S_k^2(\widehat P_k-\phi(q_k))^2\right]
=
E[S_k^2]E\!\left[(\widehat P_k-\phi(q_k))^2\right].
\]
Since \(S_k\in\{0,1\}\), \(E[S_k^2]=E[S_k]\).  The indicator integral is the probability of the pilot-light event:
\[
E[S_k]
= P^{\otimes n}\{k\in\widehat{\mathcal L}_n\}.
\]
For \(n\ge N_0\), the event \(\{k\in\widehat{\mathcal L}_n\}\) is contained in
\[
\left\{N_k^{(0)}\le 256L\right\}.
\]
The false-light condition \(p_k>B/4\), the identity \(B=4096L/m_1\), and \(m_1\le2m_0\) imply
\[
2(256L)<m_0p_k.
\]
The pilot lower-tail bound for this threshold gives
\[
E[S_k]
\le
\exp\!\left(-\frac{m_0p_k}{8}\right).
\]
Since \(p_k/B>1/4\),
\[
1+\frac{p_k}{B}\le 5\frac{p_k}{B},
\qquad
6\left(1+\frac{p_k}{B}\right)\le30\frac{p_k}{B}.
\]
Also \(6(1+p_k/B)\ge1\), so \(\log(6(1+p_k/B))\ge0\); and from \(\log x\le x-1\) for \(x>0\),
\[
\log\!\left(6\left(1+\frac{p_k}{B}\right)\right)
\le 30\frac{p_k}{B} .
\]
Together with \(M\le L\), \(B=4096L/m_1\), and \(m_1\le2m_0\), this yields
\[
2M\log\!\left(6\left(1+\frac{p_k}{B}\right)\right)
\le60L\frac{p_k}{B},
\qquad
256L\frac{p_k}{B}\le \frac{m_0p_k}{8}.
\]
Consequently
\[
-\frac{m_0p_k}{8}
+2M\log\!\left(6\left(1+\frac{p_k}{B}\right)\right)
\le -196L\frac{p_k}{B}\le -49L,
\]
and hence
\[
\exp\!\left(-\frac{m_0p_k}{8}\right)
\left\{6^M\left(1+\frac{p_k}{B}\right)^M\right\}^2
\le
\exp(-49L).
\]
Combining the preceding displays gives the single-cell bound
\[
E[e_k^2]
\le
8(\mathrm e+1)B^2\exp(-49L).
\]
% lean: falseLight_selected_error_sq

Since \(|F|\le d\), integrating the pointwise Cauchy bound and applying the single-cell estimate gives
\[
E\!\left[
\left(\sum_{k\in F}e_k\right)^2
\right]
\le
|F|\sum_{k\in F}8(\mathrm e+1)B^2\exp(-49L)
\le
d^2\,8(\mathrm e+1)B^2\exp(-49L).
\]
% lean: selectedFalseLightError_second_moment_le

Using the bounds on \(B^2\), \(\exp(-49L)\), and the rate algebra above,
\[
d^2\,8(\mathrm e+1)B^2\exp(-49L)
\le
8(\mathrm e+1)8192^2\,\frac{d^2L^2}{n^3}
\le
8(\mathrm e+1)8192^2
\left\{
\frac1n+\frac{d^2}{n^2(\log n)^2}
\right\}.
\]
% lean: selectedFalseLightError_rate
\end{proof}
